%% =====================================================================
%%  B4-T-14-02 -- what B4 contributes to "The Contrast Principle"
%%  -------------------------------------------------------------------
%%  LAYER A: APPEND-ONLY. Written once, then left alone. Material found
%%  only here sits in the `unique` slot and is protected by rule R10.
%% =====================================================================
%% @kind: source
%% @id: B4-T-14-02
%% @book: B4
%% @topic: T-14-02
%% @locator: ch. 1, pp. 13--42
%% @status: distilled
%% @unique: yes
%% @integrated: yes
%% @updated: 2026-09-19

\dossier{B4}{T-14-02}{\bookshort{B4}}{ch. 1, pp. 13--42}

\begin{distilled}
  \item A second item is perceived relative to the first; the same object reads differently depending on what preceded it, in weight, in price and in severity.
  \item The effect is demonstrated in ordinary selling sequences, where an expensive item shown first makes the intended purchase look moderate without the intended purchase changing.
  \item The principle operates on the presentation order rather than on the content of either option, so it is available to whichever party controls the sequence.
\end{distilled}

\begin{unique}
  \item Contrast as a framing mechanism distinct from scarcity and from social proof: the value shift comes from the comparison set, not from availability or from consensus.
  \item The observation that the anchor is usually an option no one intends to sell.
\end{unique}
